\documentclass[11pt]{article}
\usepackage{amsmath,amssymb,amsthm}
\newtheorem{theorem}{Theorem}
\title{A positive moment-preserving split realizes arbitrary exterior signs locally}
\author{}\date{}
\begin{document}\maketitle
\paragraph{Status.}
This is a complete written local theorem using the ordinary finite-dimensional
inverse function theorem. It is not Lean-checked. The exact checker verifies
its rational tangent identities on small examples, not the inverse function
theorem. No HeadComplexity definition is changed, and no recursive decoder
or asymptotic head saving is claimed.

For a finite measure $\mu$, write
$m_k(\mu)=\int x^k\,d\mu(x)$ and
$S_\mu(t)=\int(t-x)^{-1}\,d\mu(x)$.

\begin{theorem}[Local positive split]
Let $a<b<t_1<\cdots<t_M$, and let
$\mu=\sum_{h=1}^M w_h\delta_{p_h}$, where
$a<p_1<\cdots<p_M<b$ and $w_h>0$.
For every $\sigma\in\{-1,1\}^M$, there are positive $M$-atom
measures $\mu_+(\delta),\mu_-(\delta)$, defined for every sufficiently
small $\delta>0$, converging to $\mu$, such that
\[
 m_k(\mu_+)=m_k(\mu_-)=m_k(\mu)\quad(0\leq k<M),
\]
\[
 m_k(\mu_+)+m_k(\mu_-)=2m_k(\mu)\quad(0\leq k<2M),
\]
and
\[
 \sigma_j\bigl(S_{\mu_+}(t_j)-S_{\mu_-}(t_j)\bigr)>0
 \quad(1\leq j\leq M).
\]
The pole pairs stay disjoint and ordered around the original $p_h$.
Moreover, for every $h$, the pole of $\mu_+$ is to the right of the
corresponding pole of $\mu_-$ exactly when $\sigma_1=1$.
Thus the pair orientation is constant across all $M$ pairs, up to one
global exchange of the two children.
\end{theorem}

\begin{proof}
Put $Q(t)=\prod_h(t-p_h)$. For each index $j$ with
$\sigma_j\ne\sigma_{j+1}$, choose one root $z_j\in(t_j,t_{j+1})$.
Let $N$ be the product of the resulting linear factors, multiplied by a
nonzero constant so that $\operatorname{sign}N(t_1)=\sigma_1$.
Then $\deg N\leq M-1$, $\operatorname{sign}N(t_j)=\sigma_j$, and
$N$ has constant sign $\sigma_1$ throughout the pole interval $[a,b]$.

The partial-fraction expansion of $F=N/Q^2$ is
\[
 F(t)=\sum_h\left(\frac{\alpha_h}{t-p_h}
                  +\frac{\beta_h}{(t-p_h)^2}\right),\qquad
 \beta_h=\frac{N(p_h)}{Q'(p_h)^2},\quad
 \alpha_h=\frac{N'(p_h)}{Q'(p_h)^2}
          -\frac{N(p_h)Q''(p_h)}{Q'(p_h)^3}.              \tag{1}
\]
The formulas follow by taking the value and derivative of
$N(t)/(Q(t)/(t-p_h))^2$ at $p_h$; subtraction leaves a proper rational
function with no poles and therefore zero.

Consider the tangent variation of weights and nodes
$\dot w_h=\alpha_h$, $\dot p_h=\beta_h/w_h$.
Its Stieltjes-transform derivative is exactly $F$.
Since $F(t)=O(t^{-M-1})$ at infinity, its derivatives of moments
$m_0,\ldots,m_{M-1}$ vanish. Equivalently, expansion of (1) gives
\[
 \sum_h\bigl(\alpha_h p_h^k+k\beta_h p_h^{k-1}\bigr)=0
 \quad(0\leq k<M),                                    \tag{2}
\]
with the second summand interpreted as zero when $k=0$.

The map from the $M$ weights and $M$ ordered nodes to the first $2M$
moments has invertible Jacobian. Indeed, a tangent vector in its kernel
annihilates every polynomial of degree below $2M$ by the functional
\[
 P\longmapsto\sum_h\bigl(\dot w_hP(p_h)
                              +w_h\dot p_hP'(p_h)\bigr).
\]
Hermite interpolation independently prescribes all $2M$ values
$P(p_h),P'(p_h)$, forcing the tangent vector to vanish.
The square Jacobian is therefore invertible.

Let $v\in\mathbb R^{2M}$ be the moment derivative of the tangent
vector in (1). Its first $M$ coordinates are zero by (2).
The inverse function theorem supplies a smooth local inverse to the
moment map near $\mu$. Apply it to the two moment vectors
$m(\mu)\pm\delta v$. Their resulting measures are positive, have distinct
nodes inside $(a,b)$ for sufficiently small $\delta$, and satisfy all
the stated moment equalities exactly. Their first derivatives at zero
are the two opposite tangent vectors from (1). Consequently,
\[
 S_{\mu_+(\delta)}(t_j)-S_{\mu_-(\delta)}(t_j)
      =2\delta\frac{N(t_j)}{Q(t_j)^2}+o(\delta),
\]
which gives every strict requested sign. Similarly, the difference of
the corresponding nodes is
$2\delta\beta_h/w_h+o(\delta)$.
All $\beta_h$ have sign $\sigma_1$ and are nonzero, proving the uniform
pair orientation. The finitely many strict conditions hold at one
common sufficiently small positive scale.
\end{proof}

\paragraph{Exact tests.}
The checker \texttt{positive-moment-split-check.py} verifies (1), (2),
the query signs, and the uniform node-direction sign over every word
for $1\leq M\leq8$ on rational pole and query grids. The nonlinear
moment equalities above are proved by the local inverse argument, not
by these tests.

\paragraph{The recursive gap.}
The theorem starts both children near the \emph{same} measure $\mu$.
It does not prove that two independently prescribed descendant patterns
are both feasible near a common positive measure in their fixed-moment
fiber. Positivity, the moment budget, and the root sign construction
alone do not supply that compatibility statement. In particular, the
theorem does not replace the missing growing-capacity local compiler.
\end{document}
